\chapter{Other Graph Algorithms and Proofs}


% 3. Graph optimization (this chapter)
% Prim's - MST
%  kruska - mst
%  Topogical sort
% identifying isomorphism
%  adjacency matrices [prereq matrices and lin alg.] https://github.com/KontinuaFoundation/sequence/issues/303

% note, eliminate dijkstra and binarys search as their own chapters
This chapter is an overview of other graph algorithms, proofs, and concepts that are important to know. We will first cover adjacency matrices, then Prim's algorithm for minimum spanning trees, isomorphism, and topological sorting for directed acyclic graphs (DAGs).
\section{Adjacency Matrices}

Recall that we can represent graphs as lists of data called adjacency lists. In code, this could look like the following:
\begin{itemize}
  \item A: [B, C]
  \item B: [A, D]
  \item C: [A]
  \item D: [B]
\end{itemize}

This represents the following graph
\begin{figure}[H]
  \centering

  \begin{tikzpicture}[
      every node/.style={kontinuaNode},
      every path/.style={kontinuaEdge}
    ]

    % Nodes
    \node (A) at (0,0) {A};
    \node (B) at (1,3) {B};
    \node (C) at (3,-1) {C};
    \node (D) at (-1,2) {D};


    % Edges
    \draw (A) -- (B);
    \draw (A) -- (C);
    \draw (B) -- (D);

  \end{tikzpicture}
  \caption{A simple (undirected unweighted) graph used to make an adjacency list and adjacency matrix.}
  \label{fig:adjacency1}
\end{figure}

We could also represent the same data as a \newterm{matrix}, which we have covered in linear algebra. The \newterm{adjacency matrix} is 2D array (as compared to an adjacency list, which consists of multiple 1D arrays) where \mintinline{python}|matrix[i][j]| indicates whether an edge exists from vertex $i$ to vertex $j$. Commonly, \mintinline{python}|matrix[i][j]| $=1$ if and only if the edge from vertex $i$ to vertex $j$ exists, else \mintinline{python}|matrix[i][j]| $=0$. The above graph in Figure~\ref{fig:adjacency1} can be represented as:

$$\begin{array}{c} A=
    \begin{bmatrix}
      0 & 1 & 1 & 0 \\
      1 & 0 & 0 & 1 \\
      1 & 0 & 0 & 0 \\
      0 & 1 & 0 & 0
    \end{bmatrix}\end{array}$$

    For \emph{undirected} graphs, the adjacency matrix is symmetric, meaning that the transpose of matrix is the same matrix!
    
    $$A = A^T$$
    \begin{figure}[H]
      \centering
    
      \begin{tikzpicture}[
          every node/.style={kontinuaNode},
          every path/.style={kontinuaEdge}
        ]
    
        % Nodes
        \node (A) at (0,0) {A};
        \node (B) at (1,3) {B};
        \node (C) at (3,-1) {C};
        \node (D) at (-1,2) {D};
    
    
        % Edges
        \draw (A) -- (B);
        \draw (A) -- (C);
        \draw (B) -- (D);
        \draw (D) to [loop above, looseness=12] (D);
    
      \end{tikzpicture}
      \caption{A simple graph with an added loop on vertex D.}
      \label{fig:adjacency2}
    \end{figure}
By convention, a diagonal entry such as \mintinline{python}|matrix[k][k]|, where ($k \in \mathbb{Z}$), is non-zero only if there is a loop from a vertex to itself. The value of the entry represents the number of such loops. In this case, there are two edges from vertex (D) to itself in Figure~\ref{fig:adjacency2}, so the corresponding diagonal entry is (2).

For \emph{directed} graphs, the adjacency matrix on appears a bit different. A \newterm{directed adjacency matrix} represents a directed graph, where edges have a direction (arrows). When $A_{ij} = 1$, then an edge exists from vertex $i$ to vertex $j$.

\begin{figure}[H]
  \centering
  \begin{tikzpicture}[
      every node/.style={kontinuaNode},
      every path/.style={kontinuaEdge, directed}
    ]

    % Nodes
    \node (A) at (-1,-1) {A};
    \node (B) at (2.5,1) {B};
    \node (C) at (-1,2) {C};

    % Edges
    \draw (A) -> (B);
    \draw (A) -> (C);

  \end{tikzpicture}
  \caption{A simple directed graph used to make an adjacency matrix.}
  \label{fig:adjacency_directed}
\end{figure}

The resulting adjacency matrix is
$$\begin{array}{c} A=
\begin{bmatrix}
  0 & 1 & 1 \\
  0 & 0 & 0 \\
  0 & 0 & 0
\end{bmatrix}\end{array}$$

Notice how, now, $A \ne A^T$, because the edges are directional.

\begin{Exercise}[title={Adjacency Matrix Exercise}, label={ex:adjacency1}]



\end{Exercise}
\begin{Answer}[ref={ex:adjacency1}]



\end{Answer}
% Strengths of an adjacency matrix
\section{Prim's Algorithm and Kruskal's Algorithm for Minimum Spanning Trees}
Prim's Algorithm is minimum spanning tree
% https://github.com/KontinuaFoundation/sequence/issues/312

FIXME https://github.com/KontinuaFoundation/sequence/issues/315

\section{Travelling Salesperson Problem and Nearest Neighbour}
% https://github.com/KontinuaFoundation/sequence/issues/312
\section{Isomorphism}
\section{Topological Sort and DAGs}
%FIXME https://github.com/KontinuaFoundation/sequence/issues/342 

\section{Proofs}
%FIXME proof of leaf, subgraph
